\chapter{Results}
\label{ch:results}

\section{First setup}
With the first setup (three window zones, both doors open, 10~cm mesh, 39 settings) and transition runs, the
model has a test error of 0.11\,\% of the \cotwo{} above the outdoor level, against 0.39\,\% for persistence,
so its error is about 4 times smaller (paper's metric: 0.01\,\% against 0.05\,\%). Training with only 5, 10,
20 and 31 of the training settings gives 0.70, 0.39, 0.30 and 0.27 times the persistence error, so the
learning curve levels off. \todo{figure of a test setting from the progress report}

\section{Current setup, first dataset}
For the current setup (W1, W3, W6; 28 settings, 8~cm mesh) with the first dataset and transition runs, the
test error is 0.61\,\% against 2.74\,\% for persistence, 4.5 times smaller. The learning curve
(\tref{tab:lc}) shows that 15 training settings already give almost the same result as 20.

\begin{table}[h]
  \centering
  \caption{Learning curve, current setup, first dataset: test error of the model (5 networks) for different
  numbers of training settings. Persistence: 2.74\,\%.}
  \label{tab:lc}
  \begin{tabular}{lrrrr}
    \toprule
    Training settings & 5 & 10 & 15 & 20\\
    \midrule
    Test error & 0.76\,\% & 0.71\,\% & 0.62\,\% & 0.61\,\%\\
    \bottomrule
  \end{tabular}
\end{table}

\section{Current setup, corrected dataset}
To see what the correction of the data (\cref{ch:dataset}) means for the forecast, I trained the model twice
in exactly the same way, once on the first dataset (model A) and once on the corrected dataset (model B), first
with the direct runs. Then I tested both models on the test settings of both datasets (\tref{tab:ab}).

\begin{table}[h]
  \centering
  \caption{Test error of models trained on the first (A) and the corrected (B) dataset, direct runs, 20
  people; in brackets the paper's metric.}
  \label{tab:ab}
  \begin{tabular}{lrr}
    \toprule
     & Tested on first data & Tested on corrected data\\
    \midrule
    Model A (trained on first data) & 0.30\,\% (0.043\,\%) & 0.37\,\% (0.049\,\%)\\
    Model B (trained on corrected data) & 0.30\,\% (0.044\,\%) & 0.33\,\% (0.046\,\%)\\
    Persistence & 3.15\,\% (0.420\,\%) & 3.20\,\% (0.421\,\%)\\
    \bottomrule
  \end{tabular}
\end{table}

On the corrected data, model B is about 11\,\% more accurate than model A (0.33 against 0.37\,\%), and both are
about 10 times better than persistence. On the first data both models are equally good. The gain is in the fast
settings, where the data changed most (\tref{tab:ab_cases}); in the calm settings both models are the same.

\begin{table}[h]
  \centering
  \caption{Test error per test setting on the corrected data (persistence in brackets).}
  \label{tab:ab_cases}
  \footnotesize
  \begin{tabular}{lrrrrrrrr}
    \toprule
     & N14 & N15 & N19 & N20 & N23 & N24 & N27 & N28\\
    \midrule
    Model A & 0.23 & 0.58 & 0.12 & 0.08 & 0.25 & 1.17 & 0.23 & 0.32\\
    Model B & 0.14 & 0.59 & 0.09 & 0.08 & 0.23 & 1.16 & 0.16 & 0.21\\
    Persistence & (1.32) & (5.62) & (2.58) & (1.43) & (3.29) & (7.24) & (1.90) & (2.25)\\
    \bottomrule
  \end{tabular}
  \\[0.3em]{\small all values in \%}
\end{table}

Each model was trained only once (5 networks each), so part of the difference could be chance. That the gain
appears in exactly the fast settings and nowhere else speaks for a real effect. \todo{repeat both trainings with
other seeds}

\section{Transition runs on the corrected data}
\todo{results of models A and B trained on the transition runs (paper-like setup), to compare with 0.61\,\% vs
2.74\,\%}
